\chapter{Einleitung}
\label{chap:einleitung}

\subsubsection{Abstract}
In dieser Arbeit wird untersucht, warum sich Kryptowährungen bis 2026  nicht als normales Zahlungsmittel etabliert haben. Dazu werden drei Wirtschaftsräume verglichen: El Salvador, die Vereinigten Arabischen Emirate und die Europäische Union. Die Untersuchung basiert auf einer qualitativen Auswertung wissenschaftlicher Studien sowie offizieller Regulierungsdokumente.
% Schlussfolgerungen werden nach Abschluss der Analyse ergänzt.

\section{Relevanz und Eingrenzung des Themas}

Seit der Einführung von Bitcoin im Jahr 2009 sind über 50 Millionen weitere Kryptowährungen entstanden~\cite{yieldfund_crypto2025}. Im Euroraum besaßen im Jahr 2022 je nach Altersgruppe bis zu 12\% der Bevölkerung Krypto-Assets~\cite{BundesbankMiCAR}. Dennoch hat sich bis 2026 keine Kryptowährung als reguläres Zahlungsmittel in einer Volkswirtschaft dauerhaft etabliert~\cite{Ammous2018, Alvarez2023, BundesbankMiCAR}.
Die staatlichen Reaktionsmöglichkeiten reichen von der gesetzlichen Einführungspflicht bis hin zu strikten Verboten. Da diese Unterschiede bisher nicht systematisch verglichen wurden, ist das Thema wissenschaftlich relevant.
Die vorliegende Arbeit konzentriert sich auf drei Wirtschaftsräume: El Salvador, die Vereinigten Arabischen Emirate und die Europäische Union. Technische Aspekte von Blockchain-Protokollen sowie Länder außerhalb dieser drei Fälle werden nicht untersucht.

\section{Forschungsfrage}

\begin{tcolorbox}[colframe=lightblue, colback=lightblue2, coltitle=black, sharp corners=south, fonttitle=\bfseries, title=Normales Zahlungsmittel]
In dieser Arbeit bezeichnet der Begriff „normales Zahlungsmittel” eine Währung, die von der Bevölkerung regelmäßig für alltägliche Transaktionen genutzt wird und die als gesetzliches Zahlungsmittel staatlich anerkannt ist.
\end{tcolorbox}

Die Arbeit beantwortet folgende Forschungsfrage:
\begin{center}
    \textbf{Warum haben sich Kryptowährungen bisher nicht als normales Zahlungsmittel etabliert?}
\end{center}

Da sich die Frage nicht anhand eines einzelnen Falls beantworten lässt, sondern den Vergleich mehrerer staatlicher Strategien und ihrer Ergebnisse erfordert, handelt es sich dabei um ein komparativ-analytisches Erkenntnisproblem.

\section{Aufbau der Arbeit}
\hyperref[chap:forschungsstand]{Kapitel~\ref*{chap:forschungsstand}} stellt den Stand der Forschung dar und benennt offene Forschungslücken. Kapitel 3 beschreibt die Untersuchungsmethode und begründet die Fallauswahl. \hyperref[chap:kern]{Kapitel~\ref*{chap:kern}} analysiert die drei Regionen. \hyperref[chap:interpretation]{Kapitel~\ref*{chap:interpretation}} vergleicht die Ergebnisse und beantwortet die Forschungsfrage. Kapitel 6 fasst die Befunde zusammen und benennt Folgefragen.
